\chapter{Le modèle de données}
\label{ch:donnees}

\section{Diagramme de classes}

Deux grappes de classes cohabitent : le \textbf{contrat de collecte}, qui
représente ce que les sources ont répondu, et le \textbf{contrat du rapport}, qui
représente ce que le document affirme. Aucune classe du second n'invente un
champ qui n'existe pas dans le premier.

\begin{figure}[h]
\centering
\figbox{
\begin{tikzpicture}[
  cls/.style={draw=typhoonBlue!75,fill=typhoonBlue!4,rounded corners=1.5pt,
              inner sep=2.5pt,font=\scriptsize},
  rel/.style={draw=typhoonBlue!80,thick},
  font=\small]

  % ── Contrat de collecte ──
  \node[cls] (dr) {\begin{tabular}{@{}l@{}}
    \textbf{DiagnosticRecord}\\\hline
    adresse : string\\
    geometrie : Point (lat, lon)\\
    aleas : AléaObjet [0..13]\\
    batiment : BdnbFiche [0..1]\\
    hydro : HydroModel [0..1]\\
    provenance : BlocSource\\\hline
    validation : refuse les champs interdits
  \end{tabular}};

  \node[cls,below=0.75cm of dr] (alea) {\begin{tabular}{@{}l@{}}
    \textbf{AléaObjet}\\\hline
    code : string\\
    present : boolean\\
    \textbf{source} : string\\
    \textbf{resolution} : enum\\
    \textbf{attribution} : string\\
    intensity? : seulement si\\
    \ \ une grandeur réelle la pilote
  \end{tabular}};

  \node[cls,right=0.6cm of alea] (bdnb) {\begin{tabular}{@{}l@{}}
    \textbf{BdnbFiche}\\\hline
    batiment\_groupe\_id : string\\
    champs : 139 verbatim\\
    empreinte : géométrie\\
    hauteur, altitude, usage
  \end{tabular}};

  % ── Contrat du rapport ──
  \node[cls,right=1.4cm of dr] (rr) {\begin{tabular}{@{}l@{}}
    \textbf{ReportRequest}\\\hline
    sector : string\\
    scenario : ScenarioModel\\
    damage : DamageEstimate\\
    timestamp : string
  \end{tabular}};

  \node[cls,below=0.7cm of rr] (sc) {\begin{tabular}{@{}l@{}}
    \textbf{ScenarioModel}\\\hline
    key : extreme | moyen |\\
    \ \ frequent | faible \ ($=$ TRI)\\
    risk : string\\
    depthPeakM : float
  \end{tabular}};

  \node[cls,right=0.45cm of sc] (dmg) {\begin{tabular}{@{}l@{}}
    \textbf{DamageEstimate}\\\hline
    damageEUR : Range\\
    floodedConduitM : Range\\
    damagedBuildings : Range\\
    \dots\ (10 indicateurs)
  \end{tabular}};

  \node[cls,below=0.7cm of dmg] (asst) {\begin{tabular}{@{}l@{}}
    \textbf{AssistantResponse}\\\hline
    answer : string\\
    answered\_from\_report : bool\\
    refusal : out\_of\_scope |\\
    \ \ refused\_number | unavailable\\
    model : string
  \end{tabular}};

  % ── Relations ──
  \draw[rel] (dr) -- node[right,font=\scriptsize]{1} node[left,font=\scriptsize]{0..13} (alea);
  \draw[rel] (dr.east) -- node[above,font=\scriptsize]{0..1} (rr.west);
  \draw[rel] (alea) -- node[above,font=\scriptsize]{0..1} (bdnb);
  \draw[rel] (rr) -- node[left,font=\scriptsize]{1} (sc);
  \draw[rel] (rr) -- node[left,font=\scriptsize]{1} (dmg);
  \draw[rel,flowd] (rr.east) to[out=0,in=90] node[right,font=\scriptsize]{assistant ancré sur \emph{ce} rapport} (asst.north);

  \node[umlnote,below=0.5cm of bdnb,text width=4.3cm]
    {\textbf{Jamais dans le contrat :} \code{niveau}/\code{score},
     \code{recommandations} générées, \code{copernicus}. Le schéma les refuse.};
\end{tikzpicture}
}
\caption{Diagramme de classes du contrat. Les trois champs en gras d'\code{AléaObjet}
sont obligatoires : un aléa sans provenance n'est pas représentable.}
\label{fig:classes}
\end{figure}

\section{La provenance n'est pas une option}

Chaque objet d'aléa porte trois champs qui rendent le fait vérifiable par
l'utilisateur lui-même.

\begin{itemize}
  \item \code{source} — le producteur du fait (par exemple « Géorisques (BRGM /
    MTE) »), jamais le nom de notre module interne.
  \item \code{resolution} — la finesse réelle : \code{point-in-polygon},
    \code{building-level} ou \code{commune-level-estimate}.
  \item \code{attribution} — la mention de licence ouverte (LO 2.0), à afficher.
\end{itemize}

\begin{lstlisting}[caption={Un objet d'aléa tel qu'il sort du pipeline (extrait, valeurs réelles).}]
{
  "code": "inondation",
  "present": true,
  "source": "Georisques (BRGM / MTE)",
  "resolution": "point-in-polygon",
  "attribution": "Source Georisques -- donnees a jour au <millesime>",
  "evidence": { "in_tri": true, "classe": "extreme",
                "water_depth_m": 1.1, "origin": "tri-class" }
}
\end{lstlisting}

\begin{invariant}[title={L'argile reste une estimation communale}]
L'aléa retrait-gonflement des argiles est \textbf{toujours} publié en
\code{commune-level-estimate} tant qu'aucune source vectorielle n'est câblée. La
règle est écrite dans la constitution du projet : la changer exige un
amendement, pas un commit.
\end{invariant}

\section{La fiche bâtiment, verbatim}

La BDNB fournit 139 champs décrivant le bâtiment. Ils sont repris
\textbf{verbatim}, sans interprétation : matériaux de mur, de toiture, année de
construction, hauteur, altitude au sol, usage, empreinte. Le service n'ajoute ni
normalisation ni catégorisation.

Ce choix a un coût assumé : le contrat est volumineux, et certaines valeurs sont
vides. Il a un bénéfice direct : un expert peut vérifier chaque valeur à la
source, et aucune perte d'information n'est introduite par notre traitement. Là
où la BDNB ne répond pas, le contrat porte \code{null} — jamais une valeur
plausible inventée.

\section{Le niveau d'eau vient d'un seul endroit}

Le jumeau numérique affiche une hauteur d'eau. Cette valeur provient d'un unique
module d'entrée et n'est calculée nulle part ailleurs. En particulier, le modèle
d'impact ne doit jamais importer ce module : sans cette contrainte, une seconde
définition du niveau d'eau apparaîtrait, et deux affichages du même rapport
finiraient par diverger.

\begin{invariant}[title={Le scénario fixe le niveau, la pluie n'en fixe que la forme}]
La \emph{classe de crue TRI officielle} détermine le pic. La prévision de pluie ne
détermine que la courbe dans le temps. Hors TRI, la hauteur reste strictement
nulle. Une journée sèche produit une rampe explicitement étiquetée comme
hypothèse — elle ne fait jamais disparaître le scénario.
\end{invariant}

\section{Volumétrie}

\begin{table}[h]
\centering
\small
\begin{tabular}{@{}lrl@{}}
\toprule
\textbf{Objet} & \textbf{Volume} & \textbf{Remarque} \\
\midrule
Aléas du référentiel & 13 & toujours présents, y compris absents \\
Fiche bâtiment BDNB & 139 champs & verbatim, \code{null} si vide \\
Indicateurs d'impact & 10 & chacun en triple (valeur, basse, haute) \\
Sections du rapport & 5 & synthèse, catégories, mitigations, confiance, annexe \\
États d'absence TRI & 5 & distincts, nommés, non confondus (\S\ref{sec:etats-tri}) \\
\bottomrule
\end{tabular}
\caption{Volumétrie du contrat. Les 13 aléas sont toujours présents : c'est
l'absence qui est signifiante.}
\end{table}
